%=============================================================================
%  18_future.tex  -  Chapter 17: Future Work
%=============================================================================
\chapter{Future Work}
\label{ch:future}

Universal Analyst is a complete, working system, but its single-analyst,
single-session design leaves clear and deliberate room to grow. The roadmap
below follows naturally from the limitations stated in
Chapter~\ref{ch:security}, ordered roughly by impact.

\section{Persistence and accounts}
Sessions currently live in process memory. A persistent store --- a database for
session metadata and chat history, and durable storage for uploaded data and
generated reports --- would let analysts return to prior work, and would be the
foundation for user accounts and authentication.

\section{Multi-table and relational analysis}
Today a PDF or SQLite upload is reduced to its largest table. A natural extension
is first-class support for multiple related tables: retaining every table,
letting the analyst (or the model) express joins, and grounding answers across a
small relational schema rather than a single flat frame.

\section{Scaling the retrieval layer}
The bounded row-sampling strategy trades completeness for responsiveness at
upload time. Moving embedding to a background task with progress reporting would
allow exhaustive indexing of large datasets without blocking the request, and
approximate-nearest-neighbour tuning could keep retrieval fast at scale.

\section{Multi-tenancy and security hardening}
Authentication, authorization, per-user data isolation, and rate limiting would
turn the single-tenant tool into a shared service. Given the local-first
architecture, this can be added at the API boundary without disturbing the
grounding and validation core.

\section{Richer analytical skills}
The skill system is designed for extension. Candidate additions include
time-series forecasting, anomaly detection, automated feature engineering, and a
statistical-testing skill --- each following the established pattern of the model
proposing within a schema and deterministic code executing and measuring.

\section{Model flexibility}
The two-model strategy is configuration-driven, so supporting additional local
models --- or letting the operator choose per-task models --- is largely a matter
of configuration and light adaptation. A model-capability probe could tune
behaviour (for example, structured-output support) to whatever model is
installed.

\begin{table}[h]
\centering
\caption{Roadmap summary.}
\small
\begin{tabularx}{\linewidth}{@{}L{4.6cm} L{2.4cm} X@{}}
\toprule
\textbf{Initiative} & \textbf{Horizon} & \textbf{Primary benefit} \\
\midrule
Persistence \& accounts & Near & Return to prior work; foundation for users. \\
Multi-table analysis & Near & Real relational grounding. \\
Background indexing & Medium & Exhaustive retrieval at scale. \\
Multi-tenancy \& auth & Medium & Shared, secure deployment. \\
New analytical skills & Ongoing & Broader analytical coverage. \\
Model flexibility & Ongoing & Portability across local models. \\
\bottomrule
\end{tabularx}
\end{table}

\begin{uanote}[Growth without compromising the core]
Every item above can be added \emph{around} the grounding-and-validation core
rather than through it. The architectural decision to mediate everything behind
one backend is what keeps the roadmap additive instead of disruptive.
\end{uanote}
